323111047

\section{Introducción}

El problema central de esta actividad consiste en diseñar e implementar un algoritmo que convierta una cadena de texto de longitud variable en una huella digital o hash de tamaño fijo de 16 bytes. Para lograrlo, es necesario procesar cada letra del mensaje usando apuntadores y cambiando estados internos con operaciones lógicas y matemáticas. El reto técnico radica en lograr que la mezcla de datos sea lo suficientemente compleja para simular un algoritmo real, manejando correctamente el intercambio de bytes en la memoria y mostrando el resultado final en formato hexadecimal.

\vspace{0.4cm}

Motivación. La necesidad de realizar este trabajo surge de la importancia del hashing en la computación, ya que representa la base de la seguridad informática y de los métodos para comprobar que la información no haya sido alterada. Entender cómo funciona esta simulación de MD5 permite visualizar cómo se puede resumir información masiva en una firma única y pequeña. Además, esta práctica impulsa el uso avanzado de apuntadores y operadores de bits en el lenguaje C, los cuales son herramientas esenciales para el manejo eficiente de la memoria y el desarrollo de sistemas de software robustos.

\vspace{0.4cm}

Objetivos. El propósito principal es implementar una función de dispersión simulada que aplique conceptos de transformación y mezcla de bytes. Se busca perfeccionar el uso de apuntadores para recorrer cadenas y manipular variables por referencia mediante funciones de intercambio. Asimismo, se pretende dominar la conversión de datos binarios a una representación hexadecimal legible para el usuario, validando que cualquier cambio en la entrada produzca una salida distinta, cumpliendo así con la lógica fundamental de un algoritmo de tipo digest.